\providecommand{\Res}[2]{\qty[round-mode=figures,round-precision=3]{\fpeval{#1}}{#2}}
\providecommand{\Val}[1]{\num[round-mode=figures,round-precision=4]{\fpeval{#1}}}
\providecommand{\ResSci}[2]{\qty[round-mode=figures,round-precision=3,exponent-mode=scientific]{\fpeval{#1}}{#2}}
\subsubsection*{Question a) Fondamentales}
\begin{enumerate}[label=\arabic*)]
    \item \textbf{Variables connues et recherchées :} $L_1 = \qty{\vsi{L1}}{m}$, $L_2 = \Res{\vsi{L1}+\vsi{dL}}{m}$, $v = \qty{\v{v}}{m/s}$ ; variables recherchées : $f_1$, $f_2$
    \item \textbf{Équations pertinentes :} tuyau ouvert aux deux bouts :
    \[ f = \frac{v}{2L} \]
    \item \textbf{Isolation de la variable recherchée :} l'équation est déjà isolée.
    \item \textbf{Résolution numérique :}
    \[ f_1 = \frac{\num{\v{v}}}{2 \cdot \num{\vsi{L1}}} = \Res{\v{v}/(2*\vsi{L1})}{Hz} \qquad f_2 = \frac{\num{\v{v}}}{2 \cdot \Val{\vsi{L1}+\vsi{dL}}} = \Res{\v{v}/(2*(\vsi{L1}+\vsi{dL}))}{Hz} \]
    \item \textbf{Réponse :} $f_1 = \mathbf{\Res{\v{v}/(2*\vsi{L1})}{Hz}}$ et $f_2 = \mathbf{\Res{\v{v}/(2*(\vsi{L1}+\vsi{dL}))}{Hz}}$.
\end{enumerate}

\subsubsection*{Question b) Battements}
\[ f_{\text{batt}} = |f_1 - f_2| = \qty[round-mode=figures,round-precision=3]{\fpeval{\v{v}/(2*\vsi{L1}) - \v{v}/(2*(\vsi{L1}+\vsi{dL}))}}{Hz} \]
On entend \textbf{\Res{\v{v}/(2*\vsi{L1}) - \v{v}/(2*(\vsi{L1}+\vsi{dL}))}{Hz}} de battements (on garde plusieurs chiffres dans $f_1$ et $f_2$ avant de soustraire).

\subsubsection*{Question c) Fréquence perçue}
\[ f_{\text{moy}} = \frac{f_1 + f_2}{2} = \Res{(\v{v}/(2*\vsi{L1}) + \v{v}/(2*(\vsi{L1}+\vsi{dL})))/2}{Hz} \]
On entend une note de \textbf{\Res{(\v{v}/(2*\vsi{L1}) + \v{v}/(2*(\vsi{L1}+\vsi{dL})))/2}{Hz}} dont l'intensité varie au rythme des battements.
